\begin{table*}[t]
\centering
\small
\begin{tabular}{lrrrrrr}
\toprule
n & m & t*LP 0.05 & t*LP 0.1 & M5 binds 0.05 & M5 binds 0.1 & Thm 1(C) binds 0.1 \\
\midrule
16 & 164 & -- & -- & 4.88 & 11.2 & 26.7 \\
16 & 300 & -- & 1.25 & 6.56 & 15.6 & 36.2 \\
16 & 500 & -- & 1.5 & 8.56 & 20.6 & 46.7 \\
16 & 1,319 & 1.25 & 2.5 & 14.2 & 33.8 & 75.8 \\
16 & 10,000 & 2.5 & 4 & 39.6 & 93.2 & 209 \\
32 & 164 & -- & -- & 4.75 & 11.6 & 26.7 \\
32 & 300 & -- & 1.25 & 6.69 & 15.9 & 36.2 \\
32 & 500 & -- & 1.5 & 8.72 & 20.3 & 46.7 \\
32 & 1,319 & 1.25 & 2.5 & 14.2 & 33.9 & 75.8 \\
32 & 10,000 & 2.5 & 4 & 39.9 & 93.8 & 209 \\
64 & 164 & -- & -- & 4.86 & 11.2 & 26.7 \\
64 & 300 & -- & 1.25 & 6.56 & 15.9 & 36.2 \\
64 & 500 & -- & 1.5 & 8.61 & 20.7 & 46.7 \\
64 & 1,319 & 1.25 & 2.5 & 14.3 & 34 & 75.8 \\
64 & 10,000 & 2.5 & 4 & 40 & 94.1 & 209 \\
\bottomrule
\end{tabular}
\caption{Certified horizon against the lower bound. t*LP(eps): the largest t = k/n on the grid [1.125, 1.25, 1.5, 1.75, 2.0, 2.5, 3.0, 4.0, 6.0] whose LP certificate (95\%, sup-bias plus Hoeffding over linear estimators of the count histogram) is at most eps. M5 binds(eps): the smallest t at which the optimised two-point bound reaches eps, so that every estimator errs by at least eps with probability 1/4 under some D (scanned up to t = 1024, refined by bisection on k). Thm 1(C) binds: the same for the analytic bound. Between t*LP and the binding t it is open whether any estimator certifies eps (164 = HumanEval, 500 = MATH-500, 1,319 = GSM8K test).}
\label{tab:T10_certified_horizons}
\end{table*}
